% !TEX root = ../main.tex
% !TEX program = lualatex
\documentclass[../main.tex]{subfiles}
\IfSubfilesClassLoaded{\externaldocument{\subfix{../build/main}}}{}
% ====================
% File: 13C_cost-price-eval.tex
% ====================
\begin{document}
\IfSubfilesClassLoaded{\chapter{Cost and Price Evaluation}}{}

% === COST AND PRICE EVALUATION (EDO 3.2.4) ===
\section{Cost and Price Evaluation}

\subsection{Summary}
\begin{description}
  \item[\textbf{Fair and reasonable price.}] The \ac{ko} determines fair and reasonable price using price, cost, technical, field-pricing, and risk analyses along with sound business judgment~\autocite{edo-3-2-4-cost-price-evaluation-2025,FAR-15-404-1-Proposal-Analysis,FAR-15-404-2-Field-Pricing-Assistance}.
  \item[\textbf{Financial health signals.}] Balance sheets and income statements enable \ac{ros} and \ac{roa} calculations that inform responsibility and financial capacity decisions~\autocite{edo-3-2-4-cost-price-evaluation-2025}.
  \item[\textbf{Direct vs.\ indirect.}] Direct costs tie to a single cost objective; indirect costs (overhead and \ac{ganda}) support multiple objectives and must be allocated through defined bases~\autocite{edo-3-2-4-cost-price-evaluation-2025,FAR-31-202-Direct-Costs,FAR-31-203-Indirect-Costs}.
  \item[\textbf{Allowability tests.}] Billable costs must be reasonable, allocable, compliant with \ac{cas}/\ac{gaap}, within contract terms, and within \ac{far} 31.205 limitations~\autocite{edo-3-2-4-cost-price-evaluation-2025,FAR-31-201-2-Allowability,FAR-31-205-Selected-Costs}.
  \item[\textbf{Rate discipline.}] Indirect rates, fully burdened rates, and forward pricing rate agreements sharpen negotiation positions and protect the prenegotiation objective~\autocite{edo-3-2-4-cost-price-evaluation-2025,FAR-15-406-1-Prenegotiation-Objectives,FAR-15-407-3-FPRA}.
\end{description}

\subsection{Fair and Reasonable Price Determination}
\begin{description}
  \item[\textbf{KO responsibility.}] The \ac{ko} evaluates price reasonableness using price, cost, technical, field-pricing support, and risk analysis inputs to reach a fair and reasonable determination~\autocite{edo-3-2-4-cost-price-evaluation-2025,FAR-15-404-1-Proposal-Analysis}.
  \item[\textbf{Analysis toolkit.}] Price and cost analysis are primary techniques, supported by technical evaluations, field pricing support, and risk analysis as needed~\autocite{edo-3-2-4-cost-price-evaluation-2025,FAR-15-404-1-Proposal-Analysis,FAR-15-404-2-Field-Pricing-Assistance}.
\end{description}

\subsection{Cost or Pricing Data Requirements}
\begin{description}
  \item[\textbf{Threshold and exceptions.}] Current 10~U.S.C.\ 3702 uses a \$10M threshold for applicable prime contracts entered after 30~June~2026; the earlier inflation-adjusted regulatory threshold was \$2.5M. Check the transaction type, prime award date, actual contract terms, applicable exceptions, and current Part~215 revision rather than applying one amount to every modification or subcontract. FY~2026 NDAA Section~1826 exempts qualifying nontraditional-contractor products/services from the certified-data requirement, subject to an authorized written waiver~\autocite{USC-10-3702,USC-10-3014,DFARS-Part215-Deviation-2026,DPCAP-RFO-Index-20261020}.
  \item[\textbf{Solicitation request.}] When required, the Government must request cost or pricing data in the solicitation and typically seeks a technical analysis of labor, material, and other proposal elements~\autocite{edo-3-2-4-cost-price-evaluation-2025}.
\end{description}

\subsection{Cost and Price Analysis}
\begin{description}
  \item[\textbf{Price analysis.}] Examination of the total proposed price without evaluating separate cost elements, using competition, historical prices, or market indicators~\autocite{FAR-15-404-1-Proposal-Analysis}.
  \item[\textbf{Cost analysis.}] Evaluation of individual cost elements and profit/fee to determine how well proposed costs represent what the contract should cost under reasonable economy and efficiency~\autocite{edo-3-2-4-cost-price-evaluation-2025,FAR-15-404-1-Proposal-Analysis}.
  \item[\textbf{Cost vs.\ price lens.}] Cost analysis looks at direct labor, indirect rates, materials, and \acp{odc}; price analysis looks at total proposed price with profit/fee embedded~\autocite{edo-3-2-4-cost-price-evaluation-2025,FAR-15-404-1-Proposal-Analysis}.
\end{description}

\subsection{Field Pricing Support}
\begin{description}
  \item[\textbf{Specialist support.}] Field pricing support includes audits and specialist assessments (accountants, auditors, cost/price analysts, negotiators, engineers, production specialists, and small and disadvantaged business utilization specialists) to assist the \ac{ko} in determining reasonableness~\autocite{edo-3-2-4-cost-price-evaluation-2025,FAR-15-404-2-Field-Pricing-Assistance}.
  \item[\textbf{DCAA/DCMA thresholds.}] The coursebook audit cue remains useful: \ac{dcaa} audit assistance is normally considered for fixed-price proposals over \$10M and cost-type proposals over \$100M. Current DoD Part~215 deviation also ties field-pricing support to the certified-cost-or-pricing-data threshold and cost-type/material-weakness risk cues~\autocite{DFARS-Part215-Deviation-2026,edo-3-2-4-cost-price-evaluation-2025,FAR-15-404-2-Field-Pricing-Assistance}.
\end{description}

\subsection{Contractor Financial Health}
\begin{description}
  \item[\textbf{Balance sheet and income statement.}] Balance sheets show financial condition at a point in time; income statements summarize revenues and expenses over a period and enable profitability ratios~\autocite{edo-3-2-4-cost-price-evaluation-2025}.
  \item[\textbf{ROS and ROA.}] \ac{ros} and \ac{roa} gauge profitability and efficiency; compare ratios to industry averages and investigate sustained shortfalls~\autocite{edo-3-2-4-cost-price-evaluation-2025}.
  \item[\textbf{Responsibility risk.}] Debt, cash-on-hand, and cash-flow trends complement ratio checks to determine offeror responsibility~\autocite{edo-3-2-4-cost-price-evaluation-2025}.
\end{description}

The balance-sheet identity is assets = liabilities + equity; current liabilities are near-term obligations. An income statement moves from revenue through costs to gross, operating, and net income. The coursebook defines \ac{ros} as net income/revenue (net profit margin); an operating-profit margin uses a different numerator. State the convention and compare consistent periods and asset bases. The coursebook figures give approximately 21.4\% for 2014 and 21.8\% for 2015; profitability alone does not establish liquidity or responsibility~\autocite{edo-3-2-4-cost-price-evaluation-2025}.

\subsection{Direct and Indirect Costs}
\begin{description}
  \item[\textbf{Direct costs.}] Costs traceable to a single cost objective (direct labor, direct material, and \acp{odc}) with clear identification to a specific requirement~\autocite{edo-3-2-4-cost-price-evaluation-2025,FAR-31-202-Direct-Costs}.
  \item[\textbf{Indirect costs.}] Costs that benefit more than one objective and are allocated through overhead and \ac{ganda} pools, such as management, legal, travel, utilities, and facility security~\autocite{edo-3-2-4-cost-price-evaluation-2025,FAR-31-203-Indirect-Costs}.
\end{description}

The breakdown of the proposed price is shown in \Fig{\ref{fig:direct_indirect_costs}}.

\begin{figure}[H]
  \centering
  \includegraphics[width=\linewidth]{direct_indirect_costs.pdf}
  \caption[Direct vs.\ Indirect Costs]{Mapping direct vs.\ indirect contractor cost elements.~\srcCite{edo-3-2-4-cost-price-evaluation-2025}}
  \label{fig:direct_indirect_costs}
\end{figure}

\subsection{Indirect Cost Pools, Bases, and Rates}
\begin{description}
  \item[\textbf{Cost pools.}] Indirect cost pools group similar indirect costs (e.g., material handling, design support, building maintenance) that are allocated to cost objectives using a logical, causal-beneficial relationship~\autocite{edo-3-2-4-cost-price-evaluation-2025,FAR-31-203-Indirect-Costs}.
  \item[\textbf{Allocation bases.}] Bases include direct labor dollars or hours, material dollars, machine hours, or units produced and are used to compute indirect rates~\autocite{edo-3-2-4-cost-price-evaluation-2025,FAR-31-203-Indirect-Costs}.
\end{description}

Apply the standard relationships from the coursebook when validating proposals; they are board favorites for quick-turn computations covering indirect pools, \ac{tci}, \ac{ganda}, and profitability metrics such as \ac{ros} and \ac{roa}~\autocite{edo-3-2-4-cost-price-evaluation-2025}.
\begin{align*}
  \text{Indirect Cost Rate}        & = \frac{\text{Indirect Cost Pool}}{\text{Allocation Base}}                                               \\
  \text{TCI}                       & = \text{Direct Cost} + \text{Overhead Cost}                                                              \\
  \text{G\&A Cost}                 & = \text{G\&A Rate} \times \text{TCI}                                                                     \\
  \text{Burdened Labor Cost Rate} & = \frac{\text{Direct Labor Cost} + \text{Allocated Indirect Costs}}{\text{Direct Labor Hours}} \\
  \text{ROS}                       & = \frac{\text{Net Income}}{\text{Revenue}}                                                       \\
  \text{ROA}                       & = \frac{\text{Net Income}}{\text{Total Assets}}
\end{align*}

\note{TINA Lite is now a standing authority under 10~U.S.C.\ 3709, implemented by the July~2026 Part~215 revision. DPCAP describes its use, with approval, for actions over \$50M when the contractor has meaningful history delivering the same or similar products. It tailors certified-data requirements using historical costs; it is not a blanket certification exemption~\autocite{DPCAP-TINA-Lite-2026}.}

\paragraph{Worked allocation.} A \$2,000 overhead pool allocated over \$1,500 of direct labor produces a 133.333\% rate. Jobs with \$750, \$500, and \$250 direct labor receive \$1,000, \$666.67, and \$333.33 overhead, respectively, giving costs of \$1,750, \$1,166.67, and \$583.33 before other burdens. For a total-cost-input \ac{ganda} base, \$50,000 direct cost plus \$100,000 overhead gives \$150,000 \ac{tci}; a \$15,000 \ac{ganda} pool yields a 10\% rate and \$165,000 total cost before fee. Thus G\&A rate = G\&A pool / TCI; round final allocated dollars, not the intermediate rate~\autocite{edo-3-2-4-cost-price-evaluation-2025}.

\subsection{Cost Allowability}
\begin{description}
  \item[\textbf{Five bases.}] Reasonableness, allocability, compliance with \ac{cas}/\ac{gaap}, contract terms, and \ac{far} 31.205 limitations govern allowability~\autocite{edo-3-2-4-cost-price-evaluation-2025,FAR-31-201-2-Allowability,FAR-31-201-3-Reasonableness,FAR-31-201-4-Allocability,FAR-31-205-Selected-Costs}.
  \item[\textbf{Unallowable examples.}] Bad debts, contributions/donations, entertainment, and alcoholic beverages are unallowable~\autocite{edo-3-2-4-cost-price-evaluation-2025,FAR-31-205-3-Bad-Debts,FAR-31-205-8-Contributions,FAR-31-205-14-Entertainment,FAR-31-205-51-Alcoholic-Beverages}.
\end{description}
\note{Slides may refer to ``cash allowability''; \ac{far} terminology is \emph{cost} allowability.}

\note{FAR Part~31 remains the general cost-allowability framework. FY~2026 NDAA Section~1826 specifically exempts products/services from qualifying nontraditional defense contractors from Part~31 and specified special cost requirements unless the exemption is properly waived. Verify the 10~U.S.C.\ 3014 definition; ``new vendor'' or ``small business'' alone does not establish this status~\autocite{USC-10-3014}.}

\subsection{Current Cost Accounting Standards Thresholds}
Effective 1~October~2026, the CAS rule exempts negotiated contracts and subcontracts not exceeding \$35M and removes the old \$7.5M trigger-contract test. Full coverage and Disclosure Statement thresholds increase from \$50M to \$100M, subject to the detailed business-unit, award-history, exemption, and transition rules. These CAS tests are separate from certified cost-or-pricing-data thresholds and from \ac{evms} applicability. Check the actual contract and applicable rule before concluding that a contractor is covered~\autocite{CAS-Thresholds-2026}.

\subsection{Fully Burdened Rate}
\begin{description}
  \item[\textbf{Definition.}] A burdened labor \emph{cost} rate combines direct labor and allocated indirect costs per labor hour. A quoted billing or price rate may additionally include profit/fee; identify which convention is being used before comparing rates or estimating work~\autocite{edo-3-2-4-cost-price-evaluation-2025}.
  \item[\textbf{Example.}] A contractor with \$1.0M direct labor, \$0.5M manufacturing overhead, and \$0.2M other costs across 50,000 labor hours has a burdened cost rate of \$34 per hour, before any separately added fee. The coursebook's \$6,325/100-hour example similarly gives \$63.25 per hour. Multiply the matching labor hours by the applicable rate and avoid adding overhead or fee twice~\autocite{edo-3-2-4-cost-price-evaluation-2025}.
\end{description}

\subsection{Forward Pricing Rates}
\begin{description}
  \item[\textbf{Purpose.}] \acp{fpra} establish agreed indirect rates (typically one to three years) to speed negotiations and reduce repetitive audits by anchoring the prenegotiation objective~\autocite{edo-3-2-4-cost-price-evaluation-2025,FAR-15-407-3-FPRA}.
  \item[\textbf{Process.}] Contractors submit a \ac{fprp}; the Government reviews and negotiates rates into an \ac{fpra}, and the \ac{aco} monitors performance against the agreed rates~\autocite{edo-3-2-4-cost-price-evaluation-2025,FAR-15-407-3-FPRA}.
\end{description}

\subsection{Practitioner Steps}
\begin{enumerate}
  \item Confirm whether certified cost or pricing data are required and document the applicable exception if not; ensure the solicitation includes any required data requests~\autocite{DFARS-Part215-Deviation-2026,edo-3-2-4-cost-price-evaluation-2025}.
  \item Evaluate total price and use cost analysis when certified cost or pricing data are required, and for other data as needed; cost analysis is not merely an optional fallback in certified-data cases. Apply cost realism for cost-reimbursement proposals and integrate technical and field-pricing support~\autocite{edo-3-2-4-cost-price-evaluation-2025,FAR-Overhaul-Part15-2026}.
  \item Validate direct and indirect cost elements, confirm allocation bases, and apply the five allowability tests to each significant cost element~\autocite{FAR-31-201-2-Allowability,FAR-31-202-Direct-Costs,FAR-31-203-Indirect-Costs}.
  \item Analyze financial health using balance sheet and income statement ratios and note any responsibility concerns in the prenegotiation objective~\autocite{edo-3-2-4-cost-price-evaluation-2025}.
  \item Capture \ac{dcaa} and \ac{dcma} inputs, forward pricing rate agreements, and fully burdened rate calculations in the contract file to justify the final determination~\autocite{edo-3-2-4-cost-price-evaluation-2025,FAR-15-404-2-Field-Pricing-Assistance,FAR-15-406-1-Prenegotiation-Objectives,FAR-15-406-3-Documenting-Negotiation,FAR-15-407-3-FPRA}.
\end{enumerate}

\subsection{Checklist}
\begin{itemize}
  \item File price and cost analysis results (with technical inputs) showing how the fair-and-reasonable price determination was reached~\autocite{FAR-15-404-1-Proposal-Analysis,FAR-15-406-3-Documenting-Negotiation}.
  \item Verify each significant cost element passes the five allowability tests and that allocation bases match current practice~\autocite{FAR-31-201-2-Allowability,FAR-31-201-4-Allocability,FAR-31-203-Indirect-Costs}.
  \item Document any \ac{fpra}/\ac{fprp} used in the prenegotiation objective, or explain why legacy rates remain valid~\autocite{FAR-15-406-1-Prenegotiation-Objectives,FAR-15-407-3-FPRA}.
  \item Record \ac{ros}/\ac{roa} trends, cash flow, and \ac{dcaa}/\ac{dcma} inputs that influenced the responsibility determination~\autocite{edo-3-2-4-cost-price-evaluation-2025}.
\end{itemize}

\subsection{Quick Review}
\begin{itemize}
  \item Price analysis evaluates the total price; cost analysis evaluates individual cost elements plus profit or fee~\autocite{FAR-15-404-1-Proposal-Analysis}.
  \item The \ac{ko} makes the fair-and-reasonable determination, using technical analysis, field pricing, \ac{dcaa}/\ac{dcma} inputs, and business judgment as needed.
  \item Certified cost or pricing data depends on current threshold, award date, and exceptions; identify the threshold before answering a board scenario.
  \item Cost allowability requires reasonableness, allocability, \ac{cas}/\ac{gaap} compliance, contract-term compliance, and no \ac{far} 31.205 prohibition~\autocite{FAR-31-201-2-Allowability,FAR-31-205-Selected-Costs}.
  \item Indirect rates depend on pools and bases; check whether the allocation base has a logical causal-beneficial relationship to the cost pool~\autocite{FAR-31-203-Indirect-Costs}.
  \item \acp{fpra} and \acp{fprp} help set the prenegotiation objective; \ac{ros}, \ac{roa}, and cash-flow trends support responsibility and risk assessments~\autocite{edo-3-2-4-cost-price-evaluation-2025,FAR-15-406-1-Prenegotiation-Objectives,FAR-15-407-3-FPRA}.
\end{itemize}

%====================
% End of file
%====================
\IfSubfilesClassLoaded{
  \RenewDocumentCommand{\entryname}{}{\textbf{\color{Modern} Acronym}}
  \RenewDocumentCommand{\descriptionname}{}{\textbf{\color{Modern} Definition}}
  \printnoidxglossary[
    type=\acronymtype,
    title=Acronyms,
    style=long-booktabs
  ]}{}
\end{document}
